\chapter{El Durazno, y el régimen que lo hizo posible}

\label{cap:loteo}
Este capítulo hace dos cosas. Primero reconstruye un loteo concreto —quién es su titular,
cómo se aprobó su plano, qué quedó afuera de esa aprobación y cómo se abasteció de agua— y
después vuelve al régimen general que lo hizo posible, comparándolo con los otros
fraccionamientos del departamento. El capítulo \ref{cap:ribera} ya estableció qué sabía el Estado
sobre los arroyos que lo atraviesan y desde cuándo.

\begin{cautela}
Antes de entrar al expediente conviene decir de dónde viene el nombre. \textbf{«El Durazno o
Peñones» es una finca del departamento desde por lo menos 1932}, cuando Silvano Y.\ Murúa y su esposa pidieron su deslinde judicial junto con otras cuatro, ante el Juzgado de 1ª
Nominación en lo Civil y con perito designado. El capítulo \ref{cap:fincas} lo documenta.

Este capítulo no puede establecer si aquella finca y la matrícula sobre la que se aprobó el
loteo son la misma superficie --- median noventa años y varias particiones ---, pero el dato
importa para leer todo lo que sigue: el loteo no inventó un lugar, ocupó uno que ya tenía
nombre, dueño y límites judicialmente fijados.
\end{cautela}

\section{El titular de El Durazno}

La Resolución Delegada Nº 369 D del Ministerio de Infraestructura, del 24 de mayo de 2022, publicada en el Boletín Oficial Nº 21.240, llena con documentación oficial el casillero que la prensa no permitía llenar.

\begin{ficha}{Res. Delegada 369 D/2022}
\begin{itemize}[leftmargin=*,nosep]
\item \textbf{Expediente:} 18-37349/19-1.
\item \textbf{Oferentes de la donación:} Héctor D'Andrea\index{D'Andrea, Héctor}, en nombre propio y en representación de \textbf{Jardín Celestial S.R.L.}, CUIT 30-71157214-3.
\item \textbf{Inmueble:} matrícula Nº \textbf{6.558} del departamento La Caldera. Proyecto urbanístico ``El Durazno''.
\item \textbf{Plano:} Mensura para Unificación, Desmembramiento y Loteo Nº \textbf{1.164}, aprobado por Resolución Delegada Nº 646/21.
\item \textbf{Superficie donada:} 8.126,79 m² en tres parcelas, que la resolución consigna como ``un 3,05\% del total de la superficie''.
\item \textbf{Destino:} reserva como espacio de uso institucional de la Provincia, conforme la Ley de Catastro Nº 2.308.
\end{itemize}
\end{ficha}

Si 8.126,79 m² equivalen al 3,05 por ciento, el total ronda los 266.450 m²: unas veintiséis hectáreas y media. Es la primera magnitud documentada de un loteo caldereño en este libro. Y la cadena de expedientes queda identificada: el 18-37349/19-1 se inicia en 2019 --- el año de los aluviones ---, el plano 1.164 se aprueba en 2021, la donación se acepta en 2022.

\section{Las cintas del valle: una lectura tentativa}

\begin{cautela}
\textbf{Esta sección no documenta: observa.} Todo lo que sigue es una lectura de una imagen
exportada de un visor público, con las medidas tomadas sobre su escala gráfica.
\textbf{Ningún dato de acá está respaldado por un instrumento}, y el apartado se conserva
---en lugar de descartarse--- porque plantea una pregunta que puede contestarse con un solo
pedido, ya consignado en el apéndice \ref{cap:pedidos}. \textbf{Hasta que ese pedido vuelva
respondido, nada de este apartado debe citarse como hallazgo de este libro.}
\end{cautela}

El parcelario del visor provincial muestra, sobre el valle al sur del pueblo, un conjunto de
\textbf{parcelas largas y angostas, paralelas entre sí y perpendiculares al eje del valle}.

\lamina{lam-cintas.jpg}{lam:cintas}%
{Parcelario en cintas sobre el valle, al sur del casco urbano de La Caldera}%
{\textbf{Capa «Catastro Urbano y Rural» del visor de IDESA}, en azul, sobre base del Instituto
Geográfico Nacional y OpenStreetMap, en naranja el área edificada. La grilla es de
\textbf{un kilómetro} de lado y permite medir. Las parcelas corren perpendiculares al eje del
valle y al curso de agua.}%
{Exportación del visor \textit{mapas.salta.gob.ar}, 15 de septiembre de 2026. Se recortó el
sector; no se alteró la geometría ni la escala.}

\textbf{Medido sobre la grilla del kilómetro}, cada cinta tiene del orden de
\textbf{novecientos metros a un kilómetro de largo} y entre \textbf{veinte y treinta metros de
ancho}. El conjunto ocupa aproximadamente \textbf{un kilómetro por un kilómetro}: unas cien
hectáreas.

\textbf{Esa magnitud es lo que vuelve interesante el apartado.} El loteo El Durazno ---la primera
magnitud documentada de este capítulo--- tiene veintiséis hectáreas y media. \textbf{El conjunto
de cintas es casi cuatro veces mayor}, y su geometría no se parece a la de un loteo: los loteos
de este capítulo forman manzanas, no franjas de un kilómetro.

\begin{cautela}
\textbf{La hipótesis, formulada como tal.} Una parcela larga y angosta, perpendicular al valle,
con un extremo sobre el curso de agua, es la forma clásica de la \textbf{división agrícola de
riego}: cada fracción toca el agua. Es la geometría que el capítulo \ref{cap:fincas} sigue desde
las particiones sucesivas de las fincas coloniales.

\textbf{Si la hipótesis fuera correcta}, el mapa mostraría algo que ningún documento de este
libro muestra: que \textbf{la traza agrícola sigue dibujada en el catastro} aunque la actividad
casi haya desaparecido. El capítulo \ref{cap:tierra} documenta una caída del ochenta y ocho por
ciento de la superficie agropecuaria; el capítulo \ref{cap:trabajo} encuentra que todavía hay
ciento noventa y tres personas ocupadas en agricultura. \textbf{La forma habría quedado y el uso
se habría ido.}

\textbf{Y podría ser falsa.} Esas cintas podrían ser un fraccionamiento reciente con geometría
alargada, o una partición hereditaria del siglo XX, o varias cosas mezcladas. \textbf{Este libro
no lo sabe.}

\textbf{Lo que lo resolvería es barato}: la matrícula y el antecedente dominial de cualquiera de
esas parcelas. Con el antecedente se ve si viene de una partición antigua o de una subdivisión
reciente, y con eso el apartado se confirma o se elimina.
\end{cautela}

\begin{cautela}
Sobre el porcentaje conviene no adelantarse. La resolución acepta una donación de espacio \textit{institucional}; no dice que ése sea el total de las cesiones, y la cesión de espacio verde puede haberse instrumentado por separado. Lo que corresponde registrar es la pregunta: qué porcentaje total de la superficie de El Durazno fue cedido al dominio público, bajo qué norma y con qué destino. El capítulo \ref{cap:pdua} mostró que el proyecto de ordenanza del PDUA no fija porcentaje alguno de cesión, de modo que la respuesta habrá que buscarla en el régimen provincial de catastro o en las ordenanzas de 1982 y 1997.

Sobre el apellido del oferente cabe una advertencia expresa. La cobertura del tratamiento de la actualización del Ordenamiento Territorial de Bosques Nativos en el Senado provincial, en diciembre de 2024, consigna la intervención de un legislador de apellido D'Andrea\index{D'Andrea, Héctor}. Este libro \textbf{no afirma ni sugiere} relación alguna entre ambas personas. El apellido es frecuente y no hay en el material relevado ningún elemento que las vincule. Se deja constancia únicamente para que quien continúe verifique o descarte el punto con documentación, y no lo descubra más tarde como si fuera un hallazgo.
\end{cautela}

El catastro parcelario permite, sin embargo, una aproximación que no depende de la imagen. Los
planos de mensura del departamento se numeran en orden de registro, y cada parcela lleva el
número del último que la afectó (lámina \ref{fig:planos}). De las cincuenta y dos parcelas en
cinta que el catastro dibuja al sur del pueblo, dieciséis no tienen número de plano y doce
llevan uno de los primeros ciento cuarenta y ocho, que ya existían en 1980; sólo tres llevan
planos posteriores a 2015. \textbf{Es un indicio de partición antigua, no una prueba}: el número
dice cuándo se midió la parcela por última vez, no cuándo se dividió la tierra. La respuesta
sigue siendo el antecedente dominial.

\begin{figure}[p]
\centering
\includegraphics[width=\textwidth,height=0.84\textheight,keepaspectratio]{fig-planos-conversion.png}
\caption[La conversión, plano por plano]{\textbf{La conversión, plano por plano.} A: parcelas del departamento según el número del último plano de mensura que las afectó, agrupado en cuatro tramos; los tramos se fechan de modo aproximado con los cuatro planos que este libro documenta, recuadrados: el 148, citado por un decreto de 1980; el 962, de la prescripción de la escuela de Potrero de Gallinato (2015); el 1.078, que presentó Ayres de Quitilipi en 2020, y el 1.164 de El Durazno, aprobado en 2021. B: el corredor La Caldera--Vaqueros. Elaboración propia con el catastro parcelario de la Provincia (IDESA, consultado en septiembre de 2026) y la cartografía de radios del Censo 2022 (INDEC). Elaboración propia; dibujo de \texttt{scripts/mapa\_planos.py} del repositorio \texttt{mapas\_caldera} (apéndice \ref{cap:fuentes}).}\label{fig:planos}
\end{figure}

\section{La aprobación del plano, y lo que deja afuera}

El acto provincial que aprobó el loteo fue obtenido en su texto íntegro, y contiene el dato que ordena todo el capítulo.

\begin{ficha}{Resolución Nº 646 D del Ministerio de Infraestructura}
\begin{itemize}[leftmargin=*,nosep]
\item \textbf{Fecha:} 26 de noviembre de 2021. Boletín Oficial Nº 21.123 del 1 de diciembre. Expediente Nº \textbf{18-37.348/2019-0}. Firmante: Camacho.
\item \textbf{Objeto:} aprueba el Plano de Mensura, Desmembramiento y Loteo Nº 1.164 sobre la matrícula 6.558, \textbf{de propiedad de los Sres.\ Héctor D'Andrea\index{D'Andrea, Héctor} y Jardín Celestial S.R.L.}, correspondiente al proyecto ``El Durazno'' \textbf{en sus Etapas 1 y 2 (lámina 3/3) y Etapa 4 (lámina 2/3)}.
\item \textbf{Marco:} Título V de la Ley Nº 2.308 y su reglamentación por Decreto Nº 1.682/19.
\item \textbf{Intervenciones con foliatura:} Municipio de La Caldera, fs.\ 58, 84 y 97/101. \textbf{Secretaría de Recursos Hídricos, fs.\ 59.} Edesa S.A., fs.\ 62/67. Instituto Provincial de la Vivienda, fs.\ 60/61. Compañía Salteña de Agua y Saneamiento S.A., fs.\ 68. Dirección de Vialidad, fs.\ 106/108. Cartografía, Cálculo y Estudio de Títulos de la DGI, fs.\ 26/28 y 55; Asesoría Jurídica, fs.\ 124/125. Unidad Coordinadora de Loteos, fs.\ 137.
\item \textbf{Artículo 2º:} las matrículas individuales de cada lote se asignarán \textbf{una vez acreditada la aceptación de la donación} del espacio institucional y presentados los certificados de libre deuda del período fiscal 2021.
\end{itemize}
\end{ficha}

\subsection*{Las Etapas 3 y 5 no están}

La resolución aprueba las etapas 1, 2 y 4. \textbf{No aprueba la Etapa 3. Tampoco la 5.}

Y la sentencia del amparo consigna como afectados por los aluviones de diciembre de 2018 y enero de 2019 a los vecinos de los loteos \textbf{El Durazno I a V}. Es decir: \textbf{cinco etapas estaban vendidas y habitadas} cuando ocurrió el desastre.

Ordénense las fechas.

\textbf{Diciembre de 2018 y enero de 2019:} cinco etapas pobladas se inundan.

\textbf{2019:} se inician los expedientes 18-37.348 --- plano --- y 18-37.349 --- donación --- ante el Ministerio de Infraestructura. El mismo año de los aluviones, y el mismo año en que el Decreto 1682/19 reemplaza el régimen de 1973.

\textbf{Noviembre de 2021:} se aprueba el plano, sólo para las etapas 1, 2 y 4.

\textbf{Mayo de 2022:} se acepta la donación, condición para que se asignen las matrículas individuales.

\textbf{El loteo estuvo vendido y habitado al menos tres años antes de que la Provincia aprobara el plano de subdivisión.} Las familias que se inundaron en 2018 vivían en parcelas que aún no tenían matrícula propia, porque su asignación quedaba condicionada a un acto que ocurrió en 2022.

Y dos de las cinco etapas siguen, hasta donde alcanza este relevamiento, sin plano aprobado.

La capa catastral de Inmuebles descargada en 2026 muestra el resultado del plano 1164: \textbf{doscientas cuarenta y siete parcelas urbanas} (catastros 6714 a 6960), que suman veintisiete hectáreas con una mediana de 875 metros cuadrados, y \textbf{un remanente rural de 1.150 hectáreas} (catastro 6713). La matrícula 6558 ya no figura. Doscientos cuarenta y siete lotes con plano, contra los cuatrocientos para los que se gestionaba el agua en 2015: la diferencia es del orden de las etapas que siguen sin plano aprobado (capítulo \ref{cap:ribera}, apartado sobre la capa catastral).

\textbf{Y las dos magnitudes de este capítulo no cierran entre sí.} Si las doscientas cuarenta y siete parcelas suman veintisiete hectáreas ---sin contar calles, ochavas ni el espacio donado---, la superficie total del fraccionamiento excede las veintiséis y media que se deducen del 3,05 por ciento de la donación. Una de las dos cifras está mal: o el porcentaje del acto se calcula sobre otra base, o la medición de la capa no corresponde al mismo polígono. Lo resuelve el balance de superficies del plano 1.164. \pendiente{balance de superficies del plano de mensura, unificación, desmembramiento y loteo Nº 1.164}

\begin{ficha}{Cómo se tiene la vivienda propia --- Censo 2022}
Población en viviendas particulares propias, según la documentación de la propiedad. La Caldera: 9.187 personas; con escritura, 61,2\,\%; con \textbf{boleto de compraventa, 21,0\,\%}; con otra documentación, 9,8\,\%; sin documentación, 8,0\,\%. Provincia de Salta: con escritura, 55,3\,\%; con boleto, 10,2\,\%; con otra documentación, 14,1\,\%; sin documentación, 20,4\,\%. Departamento Capital: boleto, 6,6\,\%.

\textit{Fuente: INDEC, Censo 2022, resultados definitivos, cuadro 6.17 de condiciones habitacionales de la población, Salta; porcentajes de elaboración propia.}
\end{ficha}

\begin{cautela}
En La Caldera, uno de cada cinco habitantes de vivienda propia la tiene por boleto de compraventa: el doble del promedio provincial, más del triple que en la Capital, y la tercera proporción más alta de los veintitrés departamentos, detrás de Santa Victoria y La Candelaria. El boleto no es una irregularidad: es la forma normal de una compra cuya escritura todavía no se otorgó. Pero es exactamente la forma que adopta una compra de lote en un fraccionamiento cuyo plano no está aprobado y cuyas parcelas no tienen matrícula propia, que es lo que este capítulo documenta para El Durazno. El censo no dice qué parte de esas 1.928 personas vive en los loteos del corredor; dice que el departamento tiene una proporción de propietarios con boleto de compraventa que duplica la provincial ---la de propietarios sin escritura por cualquier causa es, en cambio, menor que la provincial---, en un territorio donde la informalidad sin papeles ---el 8 por ciento--- es menos de la mitad de la provincial. La falta de escritura, aquí, no es la del asentamiento precario: es la de la compra formal a medio terminar.
\end{cautela}

\begin{cautela}
\textbf{Una discrepancia que este libro no resuelve.} El registro de loteos que la Municipalidad elevó a Fiscalía de Estado en septiembre de 2022 (capítulo \ref{cap:amparo}) no coincide con la Resolución 646D/21: consigna como aprobadas las etapas I, III y IV, y como irregulares la II y la V. La resolución provincial aprobó la 1, la 2 y la 4. Puede tratarse de numeraciones distintas, de una aprobación posterior o de un error en una de las dos fuentes. Lo que no cambia es que en 2022 al menos dos etapas carecían de aprobación, y que la quinta figura sin ella en ambas.
\end{cautela}

\begin{cautela}
La foja 59 concentra la pregunta más productiva de todo el expediente.

Es la intervención de la \textbf{Secretaría de Recursos Hídricos} en el trámite de aprobación del plano. Ese es el mismo organismo que había aprobado, por Resolución 383 del 18 de diciembre de 2014, la determinación de la línea de ribera de ambas márgenes de los arroyos El Durazno y Guaranguay; y el mismo que, citado en el amparo, pidió el rechazo de la demanda sosteniendo que no tenía capacidad para hacer obras.

Qué dijo esa Secretaría, entre 2019 y 2021, sobre un loteo emplazado en la microcuenca que se había desbordado y cuya línea de ribera ella misma había fijado cinco años antes, es un documento de una sola foja que este libro no consultó.
\end{cautela}

\section{El texto de la Resolución 646D/21}

Se obtuvo la resolución completa, publicada en el Boletín Oficial Nº 21.123 del 1 de diciembre de 2021. Dictada el 26 de noviembre de 2021 por el Ministro de Infraestructura, en el expediente 18-37.348/2019-0, con firma de Camacho.

\begin{ficha}{Mapa de fojas del expediente 18-37.348/2019-0}
La resolución acredita, foja por foja, quién intervino:
\begin{itemize}[leftmargin=*,nosep]
\item \textbf{Municipio de La Caldera:} fs.\ 58, 84 y 97/101.
\item \textbf{Secretaría de Recursos Hídricos:} fs.\ 59.
\item \textbf{Instituto Provincial de la Vivienda:} fs.\ 60/61.
\item \textbf{Edesa S.A.:} fs.\ 62/67.
\item \textbf{Compañía Salteña de Agua y Saneamiento S.A.:} fs.\ 68.
\item \textbf{Dirección de Vialidad de Salta:} fs.\ 106/108.
\item Cartografía, Cálculo y Estudio de Títulos de la D.G.I.: fs.\ 26/28 y 55.
\item Asesoría Jurídica de la D.G.I.: fs.\ 124/125.
\item \textbf{Unidad Coordinadora de Loteos, Loteadores y Desarrolladores} (art.\ 11 del Decreto 1682/19): fs.\ 137.
\item Coordinación Legal y Técnica del Ministerio: fs.\ 128/129 y 140.
\end{itemize}
Competencia invocada: art.\ 27 de la Ley 8.171, modificada por la Ley 8.274, y Decreto 1682/19.
\end{ficha}

Dos precisiones que el texto permite y que antes eran conjetura.

\textbf{Primera: la lámina que falta.} El artículo 1 aprueba el plano 1.164 respecto de las Etapas 1 y 2 --- lámina 3/3 --- y de la Etapa 4 --- lámina 2/3. La numeración indica que el plano tiene \textbf{tres láminas}. La lámina 1/3 no se menciona en ninguna parte de la resolución. Qué contiene, y si corresponde a las etapas no aprobadas, es una pregunta de una sola foja.

\textbf{Segunda: qué condiciona la matrícula.} El artículo 2 dispone que la Dirección General de Inmuebles asignará las matrículas individuales \textbf{una vez acreditadas dos cosas}: la aceptación por la Provincia de la donación del espacio de uso institucional, y los certificados de libre deuda del período fiscal 2021.

No figura, entre las condiciones, la ejecución de obra alguna. Un apartado posterior de este capítulo muestra que el régimen anterior --- el Decreto 1.410/73, cuyo texto también se obtuvo --- condicionaba la asignación de matrícula al cumplimiento previo de las obras y servicios comprometidos, verificado por la Municipalidad. La comparación entre ambos artículos es el hallazgo que ese apartado desarrolla.

Y una lectura que corresponde hacer con cuidado, porque es fácil forzarla.

La resolución dice que los organismos de contralor intervinieron para ``garantizar que la urbanización propuesta verifique el acceso a infraestructura y servicios básicos'', y enumera seis, entre ellos Recursos Hídricos y la empresa de agua y saneamiento. La resolución no dice qué dictaminó cada uno. Un expediente puede registrar la intervención de un organismo que se expidió en contra, o con condiciones, o que se limitó a tomar conocimiento.

Por eso la foja 59 --- una sola foja de Recursos Hídricos --- sigue siendo la pregunta más productiva del expediente, y ahora se sabe exactamente dónde está. Junto a ella, la foja 68 de la empresa de agua y las fojas 58, 84 y 97/101 del Municipio, que es donde debería constar qué dijo el municipio cuyos representantes declararon después que una etapa carecía de autorización.

Una aclaración de nomenclatura, para que quien continúe este trabajo no pierda tiempo. Existe una \textbf{Resolución 383/14 de la Secretaría de Obras Públicas} del Ministerio de Economía, que aprueba el legajo técnico de un centro sanitario en el departamento Orán y se publicó en el Boletín Nº 19.346 del 21 de julio de 2014. \textbf{No tiene relación con este capítulo.} La Resolución 383 que aquí importa es de la \textbf{Secretaría de Recursos Hídricos}, del 18 de diciembre de 2014, y se conoce por el edicto publicado en el Boletín Nº 19.467 en enero de 2015. Cada organismo lleva su propia serie de numeración anual, de modo que el mismo número y el mismo año corresponden a actos distintos. Al pedirla hay que identificar el organismo, no sólo el número. Lo mismo vale para la Resolución 185/14: la Secretaría de Ambiente dictó una con ese número sobre la recepción de acumuladores eléctricos usados, publicada en el Boletín Nº 19.281 del 10 de abril de 2014, que tampoco tiene relación con la línea de ribera de la matrícula 4366 que la Secretaría de Recursos Hídricos aprobó bajo el mismo número en julio.

\subsection*{Diez días antes de la audiencia}

El 2 de febrero de 2026, el municipio de La Caldera publicó en su sitio oficial un parte sobre las tareas en curso tras el temporal de enero.

\begin{ficha}{Parte municipal, 2 de febrero de 2026 --- área Obras Públicas}
``Gracias a las gestiones realizadas por el municipio, se dispuso de \textbf{maquinaria pesada} contratada y aportada por organismos provinciales. Con estos equipos ya se ejecutaron ---y continúan desarrollándose--- tareas de \textbf{construcción de bateas} en distintos tramos del río La Caldera, además de trabajos de \textbf{encauzamiento} y la \textbf{construcción de muros de contención}.''

Se agrega que esas acciones ``se suman a las obras de encauzamiento y bateas realizadas en la zona sur de la localidad'', y que el objetivo es mejorar el escurrimiento y evitar desbordes ante nuevas lluvias.
\end{ficha}

\begin{cautela}
La fecha importa. Diez días después, el 12 de febrero, la jueza aprobó el Plan de Manejo Integral y ratificó el marco de la condena de 2021, que exigía una estrategia basada en infraestructura verde --- reforestación, relocalización de elementos obstructores --- \textbf{por sobre} la infraestructura gris, definida en la propia sentencia como obras civiles, maquinaria pesada y edificación hormigonada.

Bateas, encauzamiento, muros de contención y maquinaria pesada son, término por término, el catálogo que la sentencia subordinó. Y en noviembre de 2025, según la cobertura de prensa ---la lista de actuaciones del expediente que el capítulo \ref{cap:amparo} reconstruye no registra ese acto---, se habría ordenado a provincia y municipio reparar el daño causado por la extracción irregular de áridos, que es la actividad que altera los canales subálveos por la misma vía.

\textbf{Este libro no sostiene que esas obras incumplan la sentencia}, y hay razones serias para no afirmarlo. La condena estableció una prelación, no una prohibición: no dijo que no se hicieran obras civiles, dijo que no fueran la estrategia principal. Un temporal reciente puede justificar intervención inmediata sobre el escurrimiento, y el propio Plan aprobado puede contemplarlas. Este libro leyó el Plan, pero no pudo establecer si estas tareas concretas están comprendidas en sus medidas ni si se comunicaron al tribunal.

Lo que queda registrado es que la respuesta municipal al temporal de enero de 2026, tal como el municipio la comunica, se describe íntegramente en el vocabulario de la infraestructura gris, sin una sola mención a reforestación, a relocalización de obstructores ni a la restauración de márgenes. Y que se ejecutó mientras el expediente estaba en trámite. Contrastar ese parte con el texto del Plan es, otra vez, la comprobación que este capítulo viene señalando.
\end{cautela}

\section{Cuatrocientos lotes con agua, seis años antes del plano}

En julio de 2015 la Secretaría de Recursos Hídricos publicó, cinco días seguidos, una gestión que este libro no conocía y que reordena la cronología entera del loteo.

\begin{ficha}{Concesión de agua pública --- B.O.\ Nº 19.577 a 19.581, 10 al 16 de julio de 2015}
Expediente 0090034-211560/2014-0.

``\textbf{Magdalena Vidal\index{Vidal, Magdalena}}, \textbf{co-titular registral del inmueble Catastro Nº 4366} del Dpto.\ La Caldera, gestiona \textbf{la finalización del trámite} de concesión de agua pública subterránea de \textbf{dren horizontal}, sito en margen derecha del Río La Caldera, para \textbf{abastecimiento poblacional de 400 Lotes del Loteo `El Durazno'}, con un volumen total anual de \textbf{0,28908 hm³} comprensivos de 33 m³/hora con 24 horas de bombeo diarios, \textbf{con carácter permanente}, y para riego de espacios comunes de 3 hectáreas con una dotación de 1,575 l/s, con carácter eventual.''

Firma: Dra.\ Silvia F.\ Santamaría\index{Santamaría, Silvia F.}, Jefe del Programa Fiscalización y Control. Secretario de Recursos Hídricos: Ing.\ Alfredo Fuertes.
\end{ficha}

\begin{cautela}
Cinco cosas, y cada una toca un punto distinto del libro.

\textbf{Primera: el loteo tenía cuatrocientos lotes en 2015.} El plano de El Durazno se aprobó por Resolución 646D en \textbf{noviembre de 2021}, y sólo respecto de las etapas 1, 2 y 4. Seis años antes, el Estado provincial tramitaba el agua para \textbf{cuatrocientos lotes} de ese mismo loteo, y lo hacía bajo el rótulo de \textbf{``abastecimiento poblacional''} --- es decir, para gente que vive.

\textbf{Segunda: aparece una co-titular registral.} \textbf{Magdalena Vidal\index{Vidal, Magdalena}}, se presenta como co-titular del catastro 4366. Este libro conocía a Héctor D'Andrea\index{D'Andrea, Héctor} y a Jardín Celestial S.R.L.\ como copropietarios de la matrícula 6.558, sobre la que se aprobó el plano. La 4366 es la Finca La Caldera o Getsemaní, sobre la que recaen las líneas de ribera de sus cuatro cursos. Quiénes son sus co-titulares, en qué proporción y desde cuándo es exactamente el antecedente dominial que este libro viene pidiendo, y ahora tiene un nombre por donde empezar.

\textbf{Tercera: el organismo sabía.} La misma Secretaría que en diciembre de 2014 determinó la línea de ribera de los arroyos El Durazno y Guaranguay sobre el catastro 4366 --- sin publicar sus coordenadas --- tramitaba siete meses después, sobre el mismo catastro, el agua para cuatrocientos lotes del loteo llamado El Durazno. No son dos expedientes que se ignoran: es el mismo organismo, el mismo inmueble y el mismo nombre.

\textbf{Cuarta: la dotación.} Los 0,28908 hectómetros cúbicos anuales equivalen a 289.080 metros cúbicos, o 792 por día. Repartidos entre cuatrocientos lotes son \textbf{casi dos metros cúbicos diarios por lote}. Compárese con la concesión del Catastro 1700, de junio de 2014: veinticuatro metros cúbicos diarios para ciento veinte personas, es decir doscientos litros por persona. La de El Durazno es varias veces más generosa por unidad.

\textbf{Y quinta, que es la que exige más cuidado.} El aviso dice ``concesión de agua pública \textbf{subterránea}'' y, en el mismo párrafo, ``\textbf{con carácter permanente}''. El artículo 143 del Código de Aguas establece que ``las concesiones de uso de aguas subterráneas, salvo las que se hubieran obtenido por declaración judicial de aguas privadas, \textbf{serán eventuales}''.

Este libro \textbf{no afirma} que haya aquí una irregularidad. Un dren horizontal en la margen de un río capta agua del subálveo, y es posible que el régimen la asimile al curso superficial y no al acuífero profundo, en cuyo caso el carácter permanente sería correcto y el artículo 143 no se aplicaría. Es una cuestión técnico-jurídica que excede a este trabajo.

Lo que en efecto corresponde registrar, y es verificable, es el \textbf{contraste}: en el mismo departamento, ante el mismo organismo y con trece meses de diferencia, un abastecimiento poblacional de agua subterránea para ciento veinte personas se otorgó como \textbf{eventual}, y otro para cuatrocientos lotes se tramitó como \textbf{permanente}.
\end{cautela}

\begin{cautela}
Y la reserva que este libro se impuso --- que un dren horizontal podría asimilarse al curso superficial y quedar fuera del artículo 143 --- \textbf{no alcanza al tercer caso}.

En octubre de 2015, El Chalchal S.R.L., cotitular del \textbf{catastro 3616}, gestionó para el Loteo ``La Misión Etapa III'' la finalización de una concesión de \textbf{agua pública subterránea de pozo perforado}, para 580 habitantes en 116 lotes, ``con \textbf{carácter permanente}'': nueve metros cúbicos por hora durante veintidós horas diarias, unos 340 litros por habitante por día. El catastro 3616 es el mismo sobre el que la Resolución 142/14 había determinado, sin publicar coordenadas, la línea de ribera de los ríos Mojotoro y La Caldera (capítulo \ref{cap:ribera}). (Diez años después, en 2025, la misma firma vuelve a figurar en el Boletín como cotitular de otro inmueble del departamento, el catastro 5925, gestionando agua del arroyo Los Mártires para cien cabezas de ganado; capítulo \ref{cap:aguabaja}.) Un pozo perforado no admite discusión sobre su naturaleza: es agua subterránea en el sentido literal del artículo 143, que declara eventuales todas esas concesiones salvo las de aguas privadas declaradas judicialmente.

De modo que el cuadro queda así, y es puramente descriptivo:

\begin{center}\small
\begin{tabular}{@{}p{3.7cm}p{2.5cm}p{2.4cm}p{2.9cm}@{}}
\toprule
\textbf{Destinatario} & \textbf{Fuente} & \textbf{Fecha} & \textbf{Carácter} \\
\midrule
120 habitantes, catastro 1700 & Pozo perforado & jun.\ 2014 & \textbf{Eventual} \\
400 lotes, loteo El Durazno & Dren horizontal & jul.\ 2015 & \textbf{Permanente} \\
580 habitantes, loteo La Misión III & \textbf{Pozo perforado} & oct.\ 2015 & \textbf{Permanente} \\
\bottomrule
\end{tabular}
\end{center}

Tres gestiones de abastecimiento poblacional, ante el mismo organismo, firmadas por la misma funcionaria, en dieciséis meses. Dos de las tres son de pozo perforado y recibieron tratamientos opuestos.

\textbf{Este libro sigue sin afirmar que exista irregularidad}, y sostiene la reserva por una razón concreta: no encontró en el Código de Aguas ni en su reglamentación ---que obtuvo completos--- una norma que distinga estas situaciones, y no conoce las resoluciones finales de ninguno de los tres trámites --- los avisos publican \textbf{gestiones}, no otorgamientos. Es enteramente posible que la ley admita el carácter permanente en supuestos que estas líneas no identifican.

Lo que registra es que la diferencia existe, es verificable en el Boletín Oficial y \textbf{correlaciona con el tamaño del solicitante}: la concesión eventual es la de ciento veinte personas; las permanentes, las de cuatrocientos y quinientos ochenta. Preguntarlo es legítimo y la respuesta está en tres expedientes identificados por número.
\end{cautela}

\textbf{Un cuarto caso, del mismo año de inicio.} El expediente 0090034-88064/2015 ---iniciado en 2015, como el de La Misión III--- llega al Boletín en abril de 2019: \textbf{Valle Hermoso S.R.L.}\ «gestiona la finalización del trámite de concesión de agua pública subsuperficial (de dren horizontal) para abastecimiento poblacional del loteo Valle Hermoso», en el catastro 3.307, por 15,71 metros cúbicos por hora con dieciséis horas diarias de bombeo ---0,092 hectómetros cúbicos al año---, \textbf{con carácter permanente}. Firma la misma funcionaria, y la publicación se repite cinco veces, entre el 26 de abril y el 3 de mayo de 2019. Es otro loteo que pide agua de un dren subsuperficial, con el mismo carácter que el de El Durazno; el aviso no dice de qué curso.

Y hay una consecuencia cronológica que el capítulo debe enunciar sin adornos.

El primer aluvión ocurrió el 21 de diciembre de 2018. Para entonces, según esta gestión, el loteo El Durazno llevaba al menos tres años y medio con trámite de agua en curso para cuatrocientos lotes destinados a abastecimiento poblacional. El plano se aprobaría tres años después del aluvión.

Este libro ya había documentado que los lotes se vendieron y se habitaron antes de la aprobación. Lo nuevo es que \textbf{el Estado provincial estaba dotando de agua a esa población} mientras el fraccionamiento carecía de plano aprobado, y que el organismo que lo hacía era el mismo que había fijado la línea de ribera de los arroyos que iban a desbordarse.

\section{Quién es Jardín Celestial S.R.L.}

La sociedad aparece en este libro por tres vías documentales independientes: es copropietaria del inmueble sobre el que se aprobó el plano, figura como tercero en el expediente judicial, y su composición consta en avisos publicados en el Boletín Oficial.

\begin{ficha}{Constitución --- Boletín Oficial Nº 18.481, 1 de diciembre de 2010}
Instrumento del 13 de octubre de 2010. Publicación autorizada por el Juez de Minas y en lo Comercial de Registro. Domicilio social: Sarmiento 122, ciudad de Salta. Duración: 50 años. Capital: \$180.000 en 1.800 cuotas.

\textbf{Objeto social:} instalación, construcción y administración del funcionamiento de un \textbf{Cementerio Parque} y todo lo inherente a su explotación \textbf{en la jurisdicción del departamento La Caldera}.

\textbf{Socios:} Héctor D'Andrea\index{D'Andrea, Héctor}, 44 años, empresario, domiciliado en Finca Getsemaní, La Caldera --- 990 cuotas, 55\%; Carlos Miguel Calabró\index{Calabró, Carlos Miguel}, 19 años, estudiante, domiciliado en Los Crisantemos s/n, La Caldera --- 540 cuotas, 30\%; y Renzo Daniel Palacios, 48 años, empresario, domiciliado en la provincia de San Juan --- 270 cuotas, 15\%. Socio gerente: Palacios. Suplente: D'Andrea.
\end{ficha}

\begin{ficha}{Modificación --- Boletín Oficial Nº 19.394, 29 de septiembre de 2014}
Acta de reunión de socios del 26 de mayo de 2014. Se modifica el objeto: la sociedad ``también podrá dedicarse a la \textbf{comercialización de inmueble de terceros}, y cuando las finanzas lo permitan podrán \textbf{adquirir inmueble para su futura realización, para este último caso se requerirá un acta que contenga el acuerdo entre los socios para cada inmueble} en oportunidad que se llegará a la adquisición''. El objeto conserva además la facultad de ``construir edificios, obras de infraestructura y todas las obras necesarias''.

Se capitalizan aportes irrevocables efectuados en agosto de 2012 \textbf{para la compra de un camión Ford, dominio JT0712}, y el capital se eleva a \$545.000 en 5.450 cuotas. Participaciones resultantes: D'Andrea\index{D'Andrea, Héctor} 2.453 cuotas (45,01\%), Calabró 1.635 (30,00\%), Palacios 1.362 (24,99\%). Certifica la Dra.\ María Luján Genovese, secretaria interina del Juzgado de Minas y en lo Comercial de Registro.
\end{ficha}

Tres detalles del texto completo que conviene retener.

\textbf{El primero:} la adquisición de inmuebles quedó sujeta a \textbf{un acta de acuerdo entre los socios para cada inmueble}. No es una cláusula de estilo. Significa que cada compra de tierra por parte de la sociedad debía dejar constancia societaria propia, y que esas actas existen. Son el documento que registraría, socio por socio y fecha por fecha, cómo la firma pasó de administrar un cementerio a ser copropietaria del inmueble de un loteo.

\textbf{El segundo, que la ficha muestra y conviene leer:} el aumento de capital \textbf{reordenó las participaciones}. D'Andrea baja del 55 al 45,01 por ciento y pierde la mayoría absoluta; Palacios sube del 15 al 24,99; y Calabró queda en \textbf{30,00 por ciento exacto}, la única cifra redonda de las tres. Ocurre en el mismo acto en que el objeto social se amplía a la compra y comercialización de inmuebles. Este libro no sabe qué lo explica ---la capitalización de aportes desiguales lo produce sin necesidad de acuerdo alguno--- y lo registra porque es el único movimiento societario documentado de la firma entre su constitución y la aprobación del plano.

\textbf{El tercero:} el aumento de capital de 2014 capitaliza aportes hechos en \textbf{agosto de 2012 para comprar un camión}. Una sociedad cuyo objeto declarado era un cementerio parque capitalizó, dos años después, la compra de un vehículo de carga. Puede ser enteramente funcional a la construcción del cementerio. Se consigna porque fecha una inversión en 2012, en el período que este libro tiene menos documentado.

\begin{ficha}{Renovación del órgano de administración --- Boletín Oficial Nº 20.675, 30 de enero de 2020}
Acta de reunión ordinaria del \textbf{8 de noviembre de 2018}. Gerente Titular: Héctor D'Andrea\index{D'Andrea, Héctor}. \textbf{Gerente Suplente: Héctor Miguel Calabró\index{Calabró, Héctor Miguel}}, con el mismo número de documento que consta para él en el acta de proclamación de 2011 (capítulo \ref{cap:politica}), quien acepta el cargo y constituye domicilio especial en Sarmiento 122.
\end{ficha}

Héctor Miguel Calabró\index{Calabró, Héctor Miguel} fue \textbf{intendente de La Caldera} --- accedió en 1999 y fue reelecto en 2003 y 2007 ---, \textbf{diputado provincial} por el departamento entre 2011 y 2015, secretario de Relaciones Institucionales del Ministerio de Gobierno desde diciembre de 2015 (Decreto 89/15, B.O.\ Nº 19.684 del 17/12/2015) y \textbf{Secretario de Asuntos Municipales} de la Provincia desde noviembre de 2017 (Decreto 1634/17, B.O.\ Nº 20.155 del 30/11/2017, dictado junto con el Decreto 1608/17, que aceptó desde el 23 de noviembre de 2017 su renuncia al cargo anterior) hasta el final de la gestión Urtubey, en 2019 ---en 2018 una decisión administrativa le asignó el adicional por responsabilidad del cargo; la fecha del cese proviene de la prensa---; los dos decretos lo identifican con el mismo número de documento que el acta de proclamación de 2011 y el aviso societario de 2020, y \textbf{senador provincial} por La Caldera entre diciembre de 2021 y diciembre de 2025.

La secuencia documentada, entonces, es la siguiente:

\begin{itemize}[leftmargin=1cm,itemsep=2pt]
\item La sociedad se constituye en octubre de 2010, \textbf{durante su tercera intendencia}, con un socio de diecinueve años que comparte su apellido y su segundo nombre, suscriptor del 30\% del capital y con objeto social expresamente circunscripto \textbf{a la jurisdicción que gobierna}.
\item El objeto se amplía a la comercialización y adquisición de inmuebles en mayo de 2014, \textbf{mientras es diputado provincial por ese departamento}.
\item Él mismo ingresa al órgano de administración en noviembre de 2018, \textbf{mientras es Secretario de Asuntos Municipales de la Provincia}.
\item El plano del loteo se aprueba en noviembre de 2021, \textbf{cuando es senador provincial electo por ese departamento}, a días de asumir la banca.
\end{itemize}

\begin{cautela}
Corresponde establecer con precisión qué prueba y qué no prueba esta secuencia, porque es el punto del libro donde más fácil sería excederse.

\textbf{No prueba ilicitud.} Participar en una sociedad comercial no lo es. Que un pariente de un funcionario sea socio de una firma, tampoco. Que un funcionario provincial ---en noviembre de 2018 era Secretario de Asuntos Municipales--- acepte un cargo de gerente suplente en una empresa privada, tampoco, aunque sí plantea la pregunta por el régimen de incompatibilidades que le correspondía. Y el gerente suplente de una sociedad de responsabilidad limitada no administra: reemplaza al titular en caso de ausencia o vacancia.

\textbf{No prueba que el loteo se haya aprobado por influencia.} La aprobación provincial del plano tramitó entre 2019 y 2021, período en el cual la intendencia de La Caldera \textbf{no era ejercida por Calabró} sino por Diego Sumbay\index{Sumbay, Diego Ramiro}, en funciones desde 2019. La intervención municipal que consta en el expediente corresponde a esa gestión.

\textbf{Y hay un registro que apunta en dirección contraria y que este libro documenta.} En diciembre de 2024, siendo senador, Calabró fue uno de \textbf{sólo tres votos negativos} a la Ley provincial 8483 de revisión del Ordenamiento Territorial de Bosques Nativos. Impulsó además la prórroga de la ley de suspensión de desalojos de familias rurales. Ambos actos son restrictivos respecto del avance sobre bosque nativo y protectores respecto de la tenencia precaria.

\textbf{Lo que en cambio establece la secuencia} es una convergencia sostenida durante quince años entre el ejercicio de función pública sobre un territorio y la participación societaria en una firma que opera suelo en ese mismo territorio. Y establece que la sociedad titular del loteo que se inundó, y que es parte del expediente judicial que litiga por esa inundación, tuvo entre sus socios a alguien que comparte apellido y segundo nombre con el intendente bajo cuya gestión se constituyó.

Nada de eso equivale a una imputación, y este libro no la formula.
\end{cautela}

\begin{cautela}
Dos advertencias de método, que valen para este apartado más que para ningún otro del libro.

\textbf{Primera.} Carlos Miguel Calabró\index{Calabró, Carlos Miguel} no ejerció función pública electiva. Es una persona privada que figura en avisos comerciales publicados en el Boletín Oficial. Se lo nombra porque esos avisos son públicos y porque la sociedad de la que fue socio es el objeto central de este capítulo, no por su vínculo familiar. Toda referencia a él se limita a lo que consta en esos instrumentos.

\textbf{Segunda.} Esta línea de investigación llegó a este trabajo por un encargo formulado por un tercero interesado, con nombre puesto de antemano. Se la siguió aplicando los mismos estándares que al resto del libro: instrumentos públicos, cotejo de documentos de identidad, y consignación expresa de lo que apunta en sentido contrario. Antes de cualquier publicación corresponde ofrecer a las personas y a la sociedad nombradas la posibilidad de responder por escrito.
\end{cautela}

\section{El circuito}

El proceso por el cual una hectárea de monte de yungas se convierte en parcelas vendibles atraviesa, en el departamento La Caldera, una secuencia relativamente estable.

\begin{enumerate}
\item \textbf{Titularidad.} Parcela madre de gran superficie, de origen sucesorio o institucional.
\item \textbf{Estructuración jurídica.} Fideicomiso, sociedad o convenio con un desarrollador. Es el paso donde el titular original se desprende del riesgo sin desprenderse del beneficio.
\item \textbf{Apertura física.} Movimiento de suelo, apertura de caminos, demarcación de lotes.
\item \textbf{Habilitación ambiental.} Expediente ante la Secretaría de Ambiente provincial, dictamen técnico y audiencia pública cuando corresponde, conforme la ley 7070.
\item \textbf{Habilitación municipal.} Certificado de aptitud ambiental municipal, ordenanza del Concejo, aprobación de plano de mensura.
\item \textbf{Comercialización.} Preventa, escrituración.
\item \textbf{Provisión de servicios.} Agua, cloacas, energía, defensas hidráulicas y caminos.
\end{enumerate}

El punto de falla no es ninguno de los pasos individualmente sino el orden en que efectivamente ocurren. Los casos relevados muestran apertura física antes de habilitación --- La Falda ---, comercialización antes de provisión de servicios --- El Durazno, Solanas ---, y aprobación municipal tramitando en paralelo a sanciones ambientales provinciales todavía vigentes --- el caso denunciado por el centro vecinal de La Calderilla.

Un procedimiento cuyos pasos son correctos pero cuya secuencia es reversible no es un procedimiento: es una formalidad posterior al hecho consumado. Y el paso séptimo queda, sistemáticamente, sin resolver o trasladado al Estado.

\section{Solanas de La Caldera}

El 21 de enero de 2026 el Concejo Deliberante celebró su Sesión Extraordinaria Nº 01, convocada por el Ejecutivo para tratar cuestiones urgentes vinculadas a ordenanzas vigentes y a la situación generada por las lluvias.

\begin{ficha}{Sesión Extraordinaria Nº 01/2026}
Según cobertura de \textit{Nuevo Diario de Salta} del 23/01/2026:
\begin{itemize}[leftmargin=*,nosep]
\item \textbf{Res. 01/26:} prórroga de la Ordenanza Nº 733.
\item \textbf{Res. 02/26:} solicitud de clausura preventiva del ingreso al Loteo Cerrado \textbf{Solanas de La Caldera}.
\item \textbf{Res. 03/26:} solicitud de clausura de los campings ``Los Teros'' y ``Papi Lalo''.
\item \textbf{Res. 04/26:} pedido de partidas para obras de defensa en el barrio Santiago Apóstol y en la zona sur.
\end{itemize}
Se promulgaron además las Ordenanzas Nº 830, 831 y 832.
\end{ficha}

El pedido de clausura se funda en un hecho del 4 de enero de 2026: el ingreso al loteo cedió y provocó daños en media calzada de la Ruta Nacional 9, a la altura del kilómetro 1618, con riesgo para la seguridad vial. El medio consigna que el loteo es copropiedad de Silvina Abilés.

El emprendimiento es ofrecido al público por la inmobiliaria Cabrera Bienes Raíces, con domicilio en calle Zuviría 642 de la ciudad de Salta, que le dedica una página propia en su sitio institucional. El material gráfico fue publicado en diciembre de 2024, lo que sitúa la comercialización activa al menos desde esa fecha. La página no consigna cantidad de lotes, superficies, precios, desarrollador ni obras comprometidas: se compone íntegramente de imágenes.

La fecha es el dato que sirve. Si la comercialización estaba activa en diciembre de 2024, los artículos 144 y 145 del Código --- garantía de ejecución de obras previa y licencia municipal de comercialización --- regían plenamente sobre este emprendimiento. El pedido de acceso a la información sobre esos dos actos debe formularse en primer término respecto de Solanas, porque es el caso donde la fecha de comercialización está documentada y el hecho posterior --- el desprendimiento del ingreso sobre una ruta nacional --- también.

Que la página comercial no informe superficie, cantidad de parcelas ni obras comprometidas es además relevante bajo el artículo 149 del mismo Código, que regula acceso a la información y defensa del consumidor en urbanizaciones.

Sobre la copropiedad, el libro fija su posición aquí y la sostiene hasta el final.

El capítulo \ref{cap:tierrafiscal} documenta, con el Decreto Nº 64/23, que María Silvina Abilés fue designada Secretaria de Tierras y Bienes del Estado el 10 de diciembre de 2023; y el capítulo \ref{cap:politica}, con el registro de la Cámara de Senadores, que fue senadora provincial por el departamento La Caldera entre 2013 y 2021. Esas dos posiciones están establecidas y son legítimas.

La copropiedad del loteo, en cambio, se apoya únicamente en una frase de una nota periodística, sin respaldo registral. Con dos eslabones confirmados, el tercero \textit{parece} confirmarse por contagio; no lo está. Un decreto de designación prueba una designación, no una titularidad inmobiliaria.

Si el informe de dominio de la matrícula confirmara la copropiedad, la lectura conjunta de las tres posiciones dejaría de ser una cuestión de trayectoria y pasaría a ser una cuestión de incompatibilidad. Precisamente por eso no se publica sin ese informe. Y precisamente por eso el pedido correspondiente encabeza la sección de Inmuebles del apéndice \ref{cap:pedidos}.

\section{El régimen provincial, y lo que reconoció en 2019}

Este libro describió el circuito desde la norma municipal. Falta el nivel provincial, que es donde se aprueba el plano sin el cual ningún lote puede tener matrícula propia.

\begin{ficha}{Decreto provincial Nº 1682/19}
\begin{itemize}[leftmargin=*,nosep]
\item \textbf{Fecha:} 25 de noviembre de 2019. Boletín Oficial Nº 20.635 del 28 de noviembre. Expediente Nº 1080384-65768/2019. Firmantes: Urtubey --- Saravia --- Klix Saravia --- Simón Padrós.
\item \textbf{Objeto:} aprueba la reglamentación del Título V de la Ley Nº 2.308, que rige la creación de nuevos centros de población y la ampliación o modificación del trazado de los existentes.
\item \textbf{Deroga el Decreto Nº 1.410/1973}, que había regido el trámite de aprobación de planos de loteos urbanos y suburbanos durante \textbf{cuarenta y seis años}.
\item \textbf{Antecedente:} la Jefatura de Gabinete creó por Resolución Nº 21/2018 una Comisión ``con el objeto de tratar la problemática relacionada a la normativa vigente en materia de loteos, urbanizaciones, clubes de campo y asentamientos de esta provincia''. Sus miembros aconsejaron modificar el decreto de 1973.
\item \textbf{Anexo:} publicado como archivo vinculado en la edición web del Boletín, con texto legible y firma digital. Treinta y tres artículos en seis capítulos.
\end{itemize}
\end{ficha}

Los considerandos del decreto son, para este libro, más elocuentes que su articulado. Sostienen que la incidencia socio-económica creciente de los desarrollos inmobiliarios \textbf{``demanda una mayor injerencia estatal a los fines de tutelar la seguridad jurídica de los eventuales adquirentes de los lotes''}; que resulta necesario ajustar las exigencias del trámite \textbf{``a fin de garantizar el efectivo cumplimiento de la normativa vigente en orden a la ejecución de las obras de infraestructura de los servicios básicos a su cargo por instancia privada''}; y que la nueva norma contempla un \textbf{sistema de registro y contralor} de los proyectos de loteo y de los catastros involucrados, \textbf{``garantizando el avance de las obras de infraestructura de los servicios públicos básicos a cargo de los desarrolladores y la publicidad e información públicas necesarias a los eventuales adquirentes interesados''}.

Léase lo que ese texto admite. En noviembre de 2019, el Poder Ejecutivo provincial declaró por decreto que la seguridad jurídica de quienes compran lotes \textbf{no estaba tutelada}, que el cumplimiento de las obras de infraestructura a cargo de los desarrolladores \textbf{no estaba garantizado}, y que faltaba publicidad e información para los adquirentes.

Es exactamente el diagnóstico que este libro construyó capítulo por capítulo, formulado por la propia autoridad que debía evitarlo, y dos años antes de que la sentencia del amparo condenara a la Provincia y al municipio.

Y hay una fecha que ordena todo lo demás: el régimen que rigió hasta ese momento era el \textbf{Decreto 1.410 de 1973}. Los loteos de El Durazno, El Nogalar y Potreros de La Caldera se vendieron y poblaron bajo una reglamentación de cuarenta y seis años de antigüedad, dictada cuando Vaqueros tenía menos de dos mil habitantes y el departamento, unos pocos miles.

El Anexo con la reglamentación \textbf{sí se publicó}: la edición web del Boletín lo ofrece vinculado al decreto, como documento con texto y firma digital. Este libro lo había dado por no publicado y lo había sumado al inventario del capítulo \ref{cap:opacidad}, del que sale. Leído, responde varias de las preguntas que este capítulo dejaba abiertas.

\begin{ficha}{Anexo del Decreto 1682/19 --- lo que exige y lo que permite}
\begin{itemize}[leftmargin=*,nosep]
\item \textbf{Art.\ 3 --- inicio en el municipio:} certificado de factibilidad de localización (o dictamen del Instituto Provincial de Vivienda si el municipio no tiene estructura técnica), pre-visado del plano de conjunto y \textbf{certificado de aptitud ambiental} con intervención de la autoridad provincial.
\item \textbf{Art.\ 4 --- expediente en Inmuebles, con carácter obligatorio:} factibilidad de localización; plano de conjunto visado; \textbf{certificado de no inundabilidad condicionado, emitido por la Secretaría de Recursos Hídricos}; pre-factibilidad de agua y cloacas; factibilidad eléctrica; pre-visado de Vialidad; nota de donación del espacio institucional provincial; certificado de aptitud ambiental.
\item \textbf{Arts.\ 6 y 7 --- obras:} proyectos ejecutivos de agua, energía y acceso vial, aprobados por cada organismo; inspección y acta de recepción; certificados de provisión efectiva; y \textbf{certificado de no inundabilidad definitivo}.
\item \textbf{Art.\ 8 --- visado final municipal:} en veinte días. Si el municipio no se expide, \textbf{lo reemplaza una declaración jurada del desarrollador}.
\item \textbf{Art.\ 9 --- aprobación del plano:} certificados de efectiva provisión de agua y energía; ordenanza municipal de aceptación de la donación de calles y espacios verdes; aceptación de la donación del espacio institucional provincial; visado final o la declaración jurada; declaración sobre desagües cloacales; plano visado por el COPAIPA; \textbf{visado de la Secretaría de Recursos Hídricos}, del municipio y de Vialidad; libre deuda municipal; tasas.
\item \textbf{Art.\ 10:} el plano se aprueba por resolución ministerial, que asigna las matrículas.
\item \textbf{Art.\ 11 --- Unidad Coordinadora:} lleva el registro de propietarios, loteadores y desarrolladores; intima a las reparticiones morosas; gestiona el régimen de preventa; y \textbf{da publicidad por prensa al estado de los loteos, a las autorizaciones de preventa y a las sanciones}.
\item \textbf{Arts.\ 12 a 17 --- publicidad:} todo aviso de venta debe consignar el \textbf{número de expediente en Inmuebles y la resolución de autorización de preventa}, y las obras con su garantía; el incumplimiento viola la Ley de Defensa del Consumidor.
\item \textbf{Arts.\ 18 a 20 --- preventa:} sólo con autorización del Ministerio y visto bueno de la Unidad Coordinadora, previa presentación de proyectos aprobados ---entre ellos, un \textbf{proyecto aprobado por Recursos Hídricos que garantice la no inundabilidad y la línea de ribera, en caso de corresponder}--- y de una \textbf{garantía real o seguro de caución} por el monto de las obras.
\item \textbf{Art.\ 21 --- matrícula anticipada:} con la captación de agua terminada y el sesenta por ciento de los desagües pluviales y de las redes internas.
\item \textbf{Arts.\ 22 y 23:} intimación a seis meses y ejecución de la garantía; ante preventa sin autorización o loteo irregular, paralización de obras, multa y denuncia a Defensa del Consumidor y a la justicia penal.
\item \textbf{Art.\ 29 --- multas:} por metro cuadrado, en porcentajes de la unidad tributaria; la más alta, de 300 a 400 por ciento, por vender sin autorización de preventa.
\item \textbf{Art.\ 33 --- transición:} los loteos con trámite iniciado tienen doce meses para adecuar planos y documentación.
\end{itemize}
\end{ficha}

\begin{cautela}
Cuatro consecuencias para este capítulo.

\textbf{Primera: la no inundabilidad y la línea de ribera también estaban en el régimen de los loteos abiertos.} El anexo exige certificado de no inundabilidad para iniciar el expediente, certificado definitivo después de las obras, visado de Recursos Hídricos para aprobar el plano y, para la preventa, un proyecto de esa Secretaría que garantice la no inundabilidad y la línea de ribera cuando corresponda. La explicación que este capítulo ofrece más abajo para la diferencia con Las Vertientes ---que son regímenes distintos--- pierde buena parte de su fuerza. La Resolución 646D/21 registra que Recursos Hídricos intervino a fojas 59; no dice qué certificó. \textbf{La foja 59 es el documento que decide la cuestión.}

\textbf{Segunda: la secuencia de 1973 no se conservó; se reemplazó.} En lugar de «primero las obras, después la venta, después la matrícula», el régimen de 2019 permite vender antes, con autorización y caución, y asignar matrículas con las obras a medio hacer. La verificación municipal de 1973 se convirtió en un visado con plazo que una declaración jurada del desarrollador suple si el municipio calla.

\textbf{Tercera: el contraste con El Durazno se vuelve más preciso.} Las dos condiciones del artículo 2 de la Resolución 646D/21 ---donación y libre deuda--- son las últimas de la lista del artículo 9; las demás debían estar cumplidas en el expediente antes de aprobar el plano. Si lo estaban ---certificados de no inundabilidad y de aptitud ambiental, provisión efectiva de agua y energía--- es lo que dirían las fojas no consultadas. Y el expediente se inició en 2019, de modo que pudo tramitar bajo la transición del artículo 33.

\textbf{Cuarta: el anexo impone deberes de publicidad.} La Unidad Coordinadora debe informar por prensa el estado de los loteos, las autorizaciones de preventa y las sanciones, y todo aviso de venta debe citar el expediente y la resolución de preventa. En las fuentes que este libro consultó no apareció ninguna de esas publicaciones, y el material de venta de Solanas de La Caldera, de 2024, puede cotejarse contra el artículo 14.
\end{cautela}

\section{Las Vertientes, y lo que dejó de exigirse}

El testimonio vecinal sobre el loteo Las Vertientes --- que este capítulo recoge en su tratamiento del agua --- tiene ahora su acto administrativo. Y el acto dice más que el testimonio.

\begin{ficha}{Resolución Nº 59 D del 30 de enero de 2014 --- B.O.\ Nº 19.241}
Ministerio de Economía, Infraestructura y Servicios Públicos, expediente 18-27.372/11. Firma: Parodi.

\textbf{Art.\ 1º:} aprueba el Plano de Mensura y Subdivisión \textbf{para Club de Campo}, según la Ley 5602 y el Decreto 924/80, del inmueble de \textbf{matrícula 3989, fracción 78}, localidad de Vaqueros, denominado \textbf{``Las Vertientes''}, de propiedad del señor \textbf{Maximiliano Rodrigo Jiménez}.

\textbf{Balance de superficies:} lotes 104 ha 2.381,68 m²; recreación pasiva 18 ha 1.115,68 m²; recreación activa 2 ha 4.217,35 m²; circulaciones 18 ha 351,42 m². \textbf{Total: 143 ha 6.688,46 m²}, igual a la superficie según título; la diferencia se consigna como «--- m²».

\textit{Nota de cotejo.} Los cuatro rubros suman \textbf{142 ha 8.066,13 m²}: el total impreso los excede en \textbf{8.622,33 m²}. La discrepancia está en el propio Boletín Oficial (pág.\ 672), no en esta transcripción. O el balance omite un rubro, o una de las cifras se publicó mal; el plano del expediente 18-27.372/11 permitiría saberlo.

\textbf{Art.\ 3º:} ``las matrículas correspondientes serán asignadas tras la presentación del \textbf{Reglamento de Condominio que contemple el Servicio de Agua en forma privada}, \textbf{Resolución de la Secretaría de Recursos Hídricos de la línea de ribera} y \textbf{Certificado de Libre Deuda}''.
\end{ficha}

\begin{cautela}
Póngase esa condición al lado de la que rige para El Durazno, aprobada siete años después por el ministerio que heredó esa competencia.

\begin{center}\small
\begin{tabular}{@{}p{3.4cm}p{4.3cm}p{4.3cm}@{}}
\toprule
& \textbf{Res.\ 59D/14} --- Las Vertientes & \textbf{Res.\ 646D/21} --- El Durazno \\
\midrule
Reglamento de condominio & \textbf{Exigido} & --- \\
\textbf{Línea de ribera} & \textbf{Exigida} & \textbf{No mencionada} \\
Libre deuda & Exigido & Exigido \\
Donación de espacio institucional & --- & Exigida \\
\bottomrule
\end{tabular}
\end{center}

En enero de 2014, para un club de campo en Vaqueros, el Ministerio condicionó la asignación de matrículas a que el desarrollador presentara \textbf{la resolución de línea de ribera} de la Secretaría de Recursos Hídricos. En noviembre de 2021, para un loteo de La Caldera \textbf{atravesado por dos arroyos cuya línea de ribera ese mismo organismo había determinado en diciembre de 2014}, el mismo Ministerio condicionó las matrículas a la donación del espacio institucional y al libre deuda fiscal, sin mencionar la línea de ribera.

\textbf{Nada de esto permite afirmar que la exigencia haya sido suprimida ni que su ausencia sea deliberada.} Son dos regímenes distintos --- club de campo bajo la Ley 5602 y el Decreto 924/80 uno, loteo bajo el Título V de la Ley 2308 y el Decreto 1682/19 el otro --- y esa diferencia normativa podía explicar la diferencia de exigencias. Pero el anexo del Decreto 1682/19, leído después, exige también para los loteos abiertos la intervención de Recursos Hídricos ---certificados de no inundabilidad y, para la preventa, línea de ribera cuando corresponda---, de modo que la diferencia de regímenes explica menos de lo que parecía (véase más arriba).

Lo que queda documentado es que el requisito \textbf{existía en la práctica del organismo} y que fue formulado por escrito, en un acto público, sobre un inmueble del mismo departamento, siete años antes. Ya no se trata de preguntar si debería haberse exigido: se trata de preguntar por qué en un caso se exigió y en el otro no.
\end{cautela}

\begin{cautela}
\textbf{Y la exigencia se cumplió, con fecha.} El 7 de abril de 2014, sesenta y siete días después de la resolución que aprobó el plano, la Secretaría de Recursos Hídricos dictó la \textbf{Resolución Nº 80}, publicada en el Boletín Oficial los días 21 y 22 de abril.

\end{cautela}

\begin{ficha}{Resolución Nº 80 del 7 de abril de 2014 --- B.O.\ Nº 19.286 y 19.287}
Expediente 34-240.329/12. Aprueba la determinación de la línea de ribera realizada por la Comisión Técnica en ambas márgenes del \textbf{arroyo Quintín o Pacará}, en la zona del inmueble de \textbf{matrícula catastral Nº 3989, Fracción 78 de Finca Vaqueros}, departamento La Caldera.

\textbf{El edicto publica las coordenadas}: veinte puntos de margen izquierda y dieciséis de margen derecha, en Gauss Krüger, con sus valores completos. Firma: Dr.\ Matías J.\ Brogin\index{Brogin, Matías J.}, Asesor Legal.
\end{ficha}

\begin{cautela}

Es la misma matrícula --- 3989, fracción 78 --- y el mismo inmueble que la Resolución 59D/14 identifica como el Club de Campo Las Vertientes. \textbf{El requisito se impuso en enero y se satisfizo en abril.}

De modo que el contraste con El Durazno deja de ser una comparación entre dos redacciones y pasa a ser una comparación entre dos trayectorias completas:

\begin{center}\small
\begin{tabular}{@{}p{3.2cm}p{4.3cm}p{4.3cm}@{}}
\toprule
& \textbf{Las Vertientes} (Vaqueros) & \textbf{El Durazno} (La Caldera) \\
\midrule
Aprobación del plano & 30/01/2014 & 26/11/2021 \\
Línea de ribera & \textbf{Exigida como condición} & No mencionada \\
Resolución de ribera & \textbf{Res.\ 80/14, 67 días después} & Res.\ 383/14, siete años antes \\
Coordenadas en el edicto & \textbf{Publicadas (36 puntos)} & Remitidas a la sede \\
\bottomrule
\end{tabular}
\end{center}

En un caso el Estado pidió la línea de ribera, el desarrollador la tramitó y el resultado se publicó entero. En el otro, la línea de ribera ya existía desde hacía siete años, no figuró entre las condiciones y sus coordenadas nunca se publicaron.

\textbf{Este libro sigue sin afirmar que la diferencia sea deliberada}, pero la explicación por los regímenes distintos queda debilitada: el anexo del Decreto 1682/19 también pedía la intervención hídrica para los loteos. Pero la comparación ya no admite la duda de si el requisito era practicable. \textbf{Era practicable en sesenta y siete días}, en el mismo departamento, con el mismo organismo y el mismo firmante.
\end{cautela}

\begin{cautela}
Y hay un segundo elemento en el artículo 3º que este capítulo no puede pasar por alto. El Estado exigió un reglamento de condominio ``que contemple el \textbf{Servicio de Agua en forma privada}''.

No es una omisión ni una tolerancia: es una \textbf{condición impuesta}. Para aprobar el fraccionamiento, el Ministerio requirió que el desarrollador acreditara que resolvería el agua por su cuenta. En el mismo departamento donde los vecinos de la zona alta llevaban entonces al menos ocho años esperando la red, y donde en 2021 seguirían esperando.

La lectura que corresponde no es de reproche: un club de campo fuera del área servida no tiene otra manera de abastecerse, y exigir que lo prevea es razonable. La lectura es de sistema, y es la que este libro viene construyendo. \textbf{El Estado no ignoraba que el suelo se convertía sin red: lo daba por supuesto y lo formalizaba.} La privatización del abastecimiento no ocurrió a espaldas del acto administrativo; ocurrió dentro de él, como requisito.

Y con el artículo 143 del Código de Aguas a la vista --- todas las concesiones de agua subterránea son eventuales ---, ese servicio privado que el Ministerio exigió prever descansa, por definición legal, sobre un título precario.
\end{cautela}

\textbf{Agua.} En julio de 2021, vecinos de la zona alta de Vaqueros se manifestaron tras más de quince años de reclamos: en sus casas no hay agua, deben comprarla o esperar el camión municipal dos o tres veces por semana, mientras los emprendimientos privados de la misma zona resuelven su abastecimiento por cuenta propia. Uno de los loteos mencionados en esa cobertura, \textit{Las Vertientes}, toma su nombre y su suministro de las vertientes cercanas. En 2017 su titular había gestionado, además, una concesión de agua pública permanente de 600 metros cúbicos diarios del río Wierna para 630 familias (capítulo \ref{cap:redes}).

\subsection*{Y la zonificación municipal ya protegía ambas márgenes}

Antes de entrar en las condiciones que cada emprendimiento recibió, conviene tener presente lo
que el capítulo \ref{cap:cot} documenta: el Código de Ordenamiento Territorial declara \textbf{no
urbanizables ambas márgenes del río La Caldera} por zonificación, con independencia de que se
haya determinado o no la línea de ribera provincial.

\textbf{Pero el código no traduce a metros el ancho de esa franja: sólo existe gráficamente, en
un plano que este libro no pudo leer en escala mensurable.} De modo que este libro puede afirmar que la protección existe y
\textbf{no puede verificar si algún loteo la invade}. Es la pieza faltante más importante de
este capítulo.

\section{Dos regímenes, y lo que no alcanzan a explicar}

Este capítulo documenta que a un club de campo del departamento se le exigió resolución de
línea de ribera y a un loteo aprobado por la misma línea ministerial, siete años después, no. Hasta aquí el libro registraba la
asimetría sin poder explicarla. Una investigación académica sobre el área metropolitana de
Salta aporta la pieza que faltaba, y no es la que este libro sospechaba.

\begin{ficha}{Polliotto, Lema, Alonso Maurizzio y Reyes, \textit{Impactos ambientales asociados al crecimiento urbano}, UCASAL, 2019}
Los autores relevan los barrios cerrados del área metropolitana desde 1990 mediante imágenes
satelitales, y advierten que hay que distinguir dos figuras «ya que, si bien se trata en ambos
casos de modalidades cerradas de urbanización, \textbf{están reguladas por normativa
diferente}»:

\textbf{Clubes de campo:} \textbf{Ley 5602/80 y Decreto 924/80}.

\textbf{Urbanizaciones privadas:} \textbf{Resolución de la Junta de Catastro 25.819/97}.

Registran además que los primeros clubes de campo del área son de 1994 y 1995 --- Altos de
San Lorenzo y El Tipal ---, que la primera urbanización privada es de 2008, y que entre los
emprendimientos de los municipios vecinos figura el \textbf{Club de Campo Las Vertientes Eco
Pueblo (2014), en el municipio de Vaqueros}.

Consignan también la superficie mínima que la normativa fija para estos emprendimientos:
\textbf{doce hectáreas}.
\end{ficha}

\textbf{Son dos regímenes distintos, y el más exigente es de rango superior.}

Un club de campo se rige por una \textbf{ley provincial} y su decreto reglamentario. Una
urbanización privada, por una \textbf{resolución de la Junta de Catastro}. La diferencia de
jerarquía normativa entre una y otra es de dos escalones.

Eso explica buena parte de lo que este capítulo documentó sin poder explicar. Las Vertientes
recibió tres condiciones --- reglamento de condominio con servicio de agua privado, resolución
de línea de ribera y libre deuda --- y Chacras del Gallinato seis, incluido un decreto del
gobernador otorgando la concesión de agua. Los dos son clubes de campo, y la fuente académica
confirma que Las Vertientes lo es. El Durazno, aprobado como loteo, recibió dos.

\begin{cautela}
Conviene medir qué explica esto y qué no.

\textbf{Explica que las exigencias sean distintas.} Dos figuras jurídicas separadas, con normas
de distinto rango, producen condiciones distintas. La asimetría deja de ser una anomalía del
expediente y pasa a ser el resultado previsible de dos regímenes paralelos.

\textbf{No explica la diferencia con los loteos abiertos.} El Durazno no se rige por una resolución de la Junta de Catastro sino por una ley ---el Título V de la 2.308--- y por un decreto cuyo anexo, leído después, exige certificados hídricos, ambientales y de servicios, y caución por las obras. \textbf{Las exigencias escritas para un loteo abierto no son menores que las de un club de campo; lo que es menor es lo que la resolución de 2021 enuncia.} Lo que sí puede decirse es que la urbanización privada ---la figura con menos parcelas por hectárea que un loteo, pero cerrada--- se aprueba bajo la norma de menor jerarquía.

Y \textbf{tampoco explica el caso concreto}: el Barrio La Misión es, como El Durazno, un loteo abierto, aprobado seis años antes ---bajo el Decreto 1.410/73 y por el entonces Ministerio de Economía---, y su plano \textit{sí} descuenta la superficie del retiro de línea de ribera. La norma no lo obligaba más que a El
Durazno.
\end{cautela}

Hay un dato más de esa investigación que conviene retener, porque enmarca todo el libro. El
área metropolitana de Salta \textbf{creció un trescientos cincuenta por ciento entre 1960 y
2010}, contra un promedio del ciento setenta por ciento en los diez aglomerados argentinos que
más crecieron. Y ese crecimiento ocurrió, señalan los autores, con indicadores
socioeconómicos por debajo de la media nacional.

El corredor norte que este libro estudia es el borde de la ciudad-región que más creció del
país.

\section{Dos aprobaciones de 2015, y lo que le exigieron a cada una}

En agosto de 2015 el Ministerio de Economía aprobó, con dos días de diferencia, dos fraccionamientos del departamento. Puestos junto al de El Durazno, ordenan la escala de exigencias del período.

\begin{ficha}{Resolución Nº 579 D del 10 de agosto de 2015 --- Barrio La Misión}
Aprueba el plano de mensura y loteo de la \textbf{matrícula 4.065}, fracción 71, sección M, localidad y departamento La Caldera, \textbf{Barrio La Misión}, de propiedad del Sr.\ \textbf{Ignacio S.\ Frías García Pinto}.

Balance: superficie unificada 94 ha 4.726 m²; lotes 44 ha 3.659 m²; \textbf{fracción de reserva 15 ha 3.799 m²}; calles y ochavas 17 ha 3.515 m²; \textbf{espacios verdes 8 ha 2.690 m²}; futuro uso institucional 1 ha 9.881 m²; plaza 1 ha 2.560 m².

Y dos notas del propio plano: ``la superficie faltante se debe \textbf{al retiro ocasionado por el trazado de la línea de ribera}, en el costado sur''; y ``\textbf{la zona de RESERVA no es urbanizable}''.

\textbf{Art.\ 2º:} las matrículas individuales se asignan previa acreditación del certificado de libre deuda.
\end{ficha}

\begin{cautela}
Este plano hace, por sí solo, lo que el capítulo \ref{cap:ribera} discute durante páginas: \textbf{descuenta superficie por la línea de ribera y lo deja escrito}. Casi seis hectáreas ---la diferencia entre la superficie unificada y la suma de los rubros del balance--- que no se lotean porque el río los reclama, y quince hectáreas más declaradas expresamente no urbanizables.

No hay aquí ningún mérito extraordinario: es lo que corresponde hacer. Se lo consigna porque demuestra que \textbf{era perfectamente posible hacerlo}, en el mismo departamento, ante el mismo Ministerio y en el mismo período. Un plano puede retirarse de la ribera y decirlo en la carátula.

\textit{Y una coincidencia que conviene descartar antes de citar los dos balances: el excedente de este plano ---5 ha 8.622 m²--- y el del balance de Las Vertientes ---8.622,33 m²--- comparten los mismos cuatro dígitos. Puede ser casualidad o puede ser una cifra arrastrada de una transcripción a la otra; se resuelve cotejando las dos resoluciones contra el Boletín, rubro por rubro.}
\end{cautela}

\begin{ficha}{Resolución Nº 585 D del 12 de agosto de 2015 --- Chacras del Gallinato}
Expediente 0110018-20984/07-0. Aprueba, en el marco de la \textbf{Ley 5.602} y el Decreto 924/80, el plano de mensura, desmembramiento y subdivisión del proyecto \textbf{``Chacras del Gallinato''}, sobre la \textbf{matrícula 3.709}, de propiedad de la firma \textbf{El Gallinato S.R.L.}

Desmembramiento: \textbf{25 ha} de superficie del \textbf{arroyo Gallinato}, 20,78 ha de Ruta Provincial 11, 2,19 ha de Ruta 161 y 1.158 m² para policía. \textbf{Superficie del club de campo: 1.569 hectáreas}, de las cuales 1.475 ha de lotes, 83 ha de área recreativa pasiva y 9,97 ha de circulación \textit{(los rubros suman 1.567,97 ha; el acto no explica la diferencia)}.

\textbf{Art.\ 3º --- condiciones para las matrículas individuales:} libre deuda fiscal; \textbf{acta de recepción definitiva de EDESA} de la obra de suministro de energía; el sistema de conducción y distribución de agua, o en su caso la constitución del \textbf{consorcio de propietarios} que administre y suministre el agua; \textbf{escritura de condominio de indivisión forzosa} conforme el art.\ 9º de la Ley 5.602; reglamento de copropiedad que prevea expresamente \textbf{la prohibición de subdividir los lotes en fracciones inferiores a cinco hectáreas}; y \textbf{decreto del Poder Ejecutivo de concesión de uso de agua pública subterránea}.
\end{ficha}

\begin{cautela}
Mil quinientas sesenta y nueve hectáreas. Es más de diez veces el club de campo Las Vertientes y más de dieciséis veces el Barrio La Misión. Es, con diferencia, el fraccionamiento más grande que este libro documenta, y su expediente lleva número de \textbf{2007}: ocho años de trámite.

Y sus condiciones son las más exigentes de la serie. Se le pide energía recibida por la distribuidora, sistema de agua o consorcio que la administre, condominio de indivisión forzosa, un reglamento que impida subdividir por debajo de cinco hectáreas, y \textbf{un decreto del gobernador otorgando la concesión de agua}. No una constancia: un decreto.

Ordénese entonces la escala, que es el punto de este apartado:

\begin{center}\small
\begin{tabular}{@{}p{3.5cm}p{3.0cm}p{5.6cm}@{}}
\toprule
\textbf{Fraccionamiento} & \textbf{Año} & \textbf{Condiciones para las matrículas} \\
\midrule
Las Vertientes & 2014 & Reglamento de condominio con agua privada, \textbf{resolución de línea de ribera}, libre deuda \\
Barrio La Misión & 2015 & Libre deuda. \textit{El plano ya descuenta la ribera y declara reserva no urbanizable} \\
Chacras del Gallinato & 2015 & Libre deuda, recepción de EDESA, sistema de agua o consorcio, condominio de indivisión forzosa, reglamento con mínimo de 5 ha, \textbf{decreto de concesión de agua} \\
\textbf{El Durazno} & \textbf{2021} & \textbf{Libre deuda y donación del espacio institucional} \\
\bottomrule
\end{tabular}
\end{center}

El fraccionamiento cuyos arroyos se desbordaron casi tres años antes de su aprobación es el que menos condiciones recibió. \textbf{Este libro no sostiene que eso obedezca a una decisión deliberada}, y mantiene la explicación disponible: los regímenes son distintos --- club de campo por Ley 5.602, loteo por el Título V de la Ley 2.308 ---, y las exigencias varían legítimamente con la figura, aunque el anexo del Decreto 1682/19 muestra que la figura de loteo también exigía la intervención hídrica.

Y el caso del Barrio La Misión, ya señalado, impide atribuir la diferencia enteramente a la figura jurídica.
\end{cautela}

\begin{ficha}{La línea de ribera de Chacras del Gallinato, 2011--2012}
\textbf{Comisión}: Resolución 396 del 27/08/2011, rectificada por la 418 del 02/09/2011 (B.O.\ Nº 18.673 y 18.674, 20 y 21 de septiembre de 2011). Ambas márgenes del arroyo Gallinato, «Lote Chacras del Gallinato», catastro 3709. Integran el ingeniero civil Marcelo Ricardo Toigo y el geólogo José A.\ Medina; trabajos desde el 12 de septiembre de 2011. Firma el asesor legal Matías J.\ Brogin\index{Brogin, Matías J.}.

\textbf{Determinación}: Resolución 011 del 17/02/2012 y su rectificatoria 020 del 05/03/2012 (B.O.\ Nº 18.826 y 18.827, 11 y 14 de mayo de 2012; también en \textit{El Tribuno}). Ambas márgenes del arroyo El Gallinato, «en el tramo que abarca la finca Mojotoro, Santa Gertrudis y Potrero de Gallinato». Las coordenadas Gauss-Krüger se ponen «a consideración de los interesados» y el instrumento puede «consultarse» en la Secretaría: \textbf{el edicto no las publica}. Firma el jefe del Programa Jurídico, Rafael Ángel Figueroa\index{Figueroa, Rafael Ángel}; edicto del 23 de abril de 2012.
\end{ficha}

\begin{cautela}
La línea de ribera del club de campo quedó determinada en 2012, tres años antes de que Inmuebles aprobara el plano, y de ella sale presumiblemente la cifra que el plano desmembra: las 25 hectáreas del arroyo Gallinato. Por eso no figura entre las condiciones de la Resolución 585 D/15: ya estaba cumplida.  La secuencia ---línea primero, plano después--- es la misma que siguió El Durazno, con su determinación de 2014 y su plano de 2021; lo que este capítulo compara no es el orden de los actos sino lo que se construyó entre uno y otro.

Dos rasgos del trámite, que el capítulo \ref{cap:ribera} retoma: cada etapa necesitó una resolución rectificatoria, y el edicto de 2012 ya remitía las coordenadas a la sede, dos años antes de las Resoluciones 142/14 y 383/14.
\end{cautela}

\begin{cautela}
Un apunte sobre nombres, registrado sin conclusión. La firma propietaria de las mil quinientas sesenta y nueve hectáreas se llama \textbf{El Gallinato S.R.L.}, y el desmembramiento incluye el \textbf{arroyo Gallinato}. El capítulo \ref{cap:tierrafiscal} documenta que en 2015 la Provincia adquirió por prescripción administrativa una fracción de la \textbf{Finca Potrero de Gallinato} para una escuela; y el artículo 6º del PDUA, del capítulo \ref{cap:pdua}, enumera una \textbf{Finca El Gallinato}.

Cuatro apariciones del mismo topónimo en cuatro actos de distinta naturaleza y distinta época. Puede tratarse de una única unidad territorial histórica fraccionada a lo largo del tiempo, o de un nombre de paraje que se repite. \textbf{Sólo el antecedente dominial lo resuelve}, y es el mismo que este libro viene pidiendo desde su primer capítulo.
\end{cautela}

\subsection*{Hubo audiencia pública para una urbanización, y fue en el departamento}

Existe un procedimiento ambiental para las urbanizaciones, y en 2014 se aplicó a un proyecto del corredor.

\begin{ficha}{Convocatoria a Audiencia Pública --- B.O.\ Nº 19.385 a 19.387, septiembre de 2014}
El Secretario de Ambiente de la Provincia convoca a Audiencia Pública ``para que quienes tengan interés o derecho puedan expresarse sobre el \textbf{impacto ambiental y social} del proyecto denominado \textbf{`Urbanización Cerro de Buena Vista'}, en inmueble \textbf{matrículas Nº 1713 y 1905} del Dpto.\ La Caldera'', expediente 0090227-97704/2013-0, presentado por los señores \textbf{Sebastián, Nicolás, Francisco y Carlos Marcos Gianzanti}, conforme el \textbf{artículo 49 de la Ley 7070}.

\textbf{Fecha:} viernes 10 de octubre de 2014, 9 de la mañana. \textbf{Lugar: Centro Deportivo Municipal de Vaqueros}, calle Los Sauces s/n esquina Los Ceibos.

El expediente y el Estudio de Impacto Ambiental y Social pueden verse en la capital y también \textbf{en la Municipalidad de Vaqueros}, avenida San Martín 1450, de lunes a viernes de 8 a 12.

Instructor de la etapa preparatoria: Ing.\ Ricardo Osvaldo Barbarán. Presidenta de la audiencia: Dra.\ Gloria Liliana Manresa.
\end{ficha}

\begin{cautela}
Conviene separar tres asuntos, y la tercera es la que este capítulo necesitaba.

\textbf{Primera: el mecanismo existe y se usa.} El artículo 49 de la Ley 7070 prevé audiencia pública sobre el impacto ambiental y social de los proyectos, y en 2014 se convocó una para una urbanización del departamento La Caldera. No es un instrumento teórico ni reservado a la minería o a la obra grande: se aplicó a un fraccionamiento privado.

\textbf{Segunda: se hizo donde vive la gente afectada.} La audiencia se celebró en el Centro Deportivo Municipal de Vaqueros, no en la capital, y el expediente completo con su estudio de impacto quedó disponible en la propia Municipalidad. Compárese con las coordenadas de línea de ribera que quedaron ``a disposición de los interesados'' en Avenida Bolivia 4650. El Estado sabe acercar los expedientes cuando decide hacerlo.

\textbf{Tercera, y es la pregunta que abre:} este libro no encontró ninguna convocatoria equivalente para el loteo El Durazno. Ni audiencia pública, ni estudio de impacto ambiental y social publicado, ni constancia de que el artículo 49 se haya activado.

Eso \textbf{no prueba} que no haya existido: el relevamiento del Boletín Oficial de este libro no es exhaustivo para todos los años, y un estudio de impacto puede tramitarse sin llegar a audiencia si la autoridad ambiental categoriza el proyecto de otro modo. Lo afirmable es que \textbf{el procedimiento estaba disponible, era conocido y se aplicaba en el departamento} en el mismo período en que los loteos del corredor se comercializaban.
\end{cautela}

\begin{cautela}
Y la finca tiene su deslinde publicado, con un siglo de anticipación.

El Boletín Oficial Nº 843, del 28 de enero de 1921, publica el edicto por el cual el juez Daniel Etcheverry ordena el \textbf{deslinde, mensura y amojonamiento de las fincas denominadas ``Cerro de Buena Vista'' y ``Sausal o Curuzú''}, sitas en el \textbf{partido de Vaqueros, departamento La Caldera}, de propiedad de don \textbf{Francisco S.\ Urquiza\index{Urquiza, Francisco S.}}. Agrimensor designado: Héctor Chiostri. Las operaciones debían comenzar el 20 de enero de 1921 \textit{(fecha anterior a la de esa edición: el edicto corría por treinta días y pudo publicarse antes en otros números)}.

Noventa y tres años después, la Secretaría de Ambiente convocaba a audiencia pública por el impacto ambiental y social de la \textbf{``Urbanización Cerro de Buena Vista''}, sobre las matrículas 1713 y 1905 del mismo departamento.

Es la cadena más completa que este libro puede exhibir de punta a punta: \textbf{una finca deslindada judicialmente en 1921, con agrimensor designado y edictos por treinta días, que en 2014 llega a audiencia pública ambiental como urbanización}. En el medio, dos actos de línea de ribera sobre una de esas matrículas, la 1905: uno de 2013 sobre el arroyo Los Nogales (Resolución 155) y otro de 2014 sobre el río Vaqueros. Y al final, el agua: en agosto de 2015 uno de los cotitulares de las matrículas 1713 y 1905 gestionó la finalización de la concesión para el \textbf{«Loteo Cerro de Buena Vista», de 51 lotes}: 63,75 metros cúbicos diarios de un dren horizontal en el río Vaqueros, con carácter permanente, y riego eventual para 3,8 hectáreas de espacios comunes (capítulo \ref{cap:aguabaja}). Y en julio de 2016 el municipio de Vaqueros aceptó, por la Ordenanza 495/16 ---promulgada el 7 de julio por el Decreto 147/16---, la donación de las calles y pasajes del «Loteo Cerros de Buena Vista» sobre las matrículas 1713 y 1905: 2 hectáreas 4.333,21 m\textsuperscript{2} de calles, 215,64 m\textsuperscript{2} de ochavas y \textbf{7.102,65 m\textsuperscript{2} de espacios verdes}, según los planos presentados ante Inmuebles (expediente 18-28358; expediente municipal 990/12). Los espacios verdes públicos donados son menos de una quinta parte de las 3,8 hectáreas de espacios comunes que el edicto de agua preveía regar; la diferencia puede ser espacio común privado del loteo, y el plano lo diría.

Este libro \textbf{no afirma} que la finca de 1921 y las matrículas de 2014 sean exactamente la misma superficie: un deslinde de hace un siglo puede haberse subdividido, vendido en partes o cambiado de límites muchas veces. Lo que afirma es que \textbf{el nombre y el lugar coinciden}, que el expediente de deslinde de 1921 existe y está identificado, y que ése es el punto de partida natural para reconstruir la cadena de títulos de ese inmueble --- algo que, para El Durazno, este libro todavía no puede hacer.
\end{cautela}

Y hay un detalle que conecta con el capítulo \ref{cap:ribera}. La urbanización propuesta se asienta, entre otras, sobre la \textbf{matrícula 1905}. Es la misma matrícula sobre la cual, cinco meses antes, la Resolución 96/14 había designado la comisión técnica para determinar la línea de ribera del río Vaqueros, y sobre la cual la Resolución 209/14 aprobó esa determinación en julio, \textbf{con las coordenadas publicadas}.

De modo que para el Cerro de Buena Vista la secuencia se cumplió completa y a la vista: comisión técnica en abril, determinación publicada con coordenadas en julio, audiencia pública ambiental en octubre, con el expediente disponible en el municipio. Es exactamente el circuito que este libro reclama, y ocurrió.

Ocurrió en Vaqueros, en 2014, para un proyecto presentado por cuatro titulares de apellido Gianzanti. Qué pasó con ese mismo circuito en el otro municipio y para el otro loteo es la pregunta que atraviesa este capítulo.

\subsection*{Quién era, en 2021, la Unidad Coordinadora}

El Decreto 1682/19 creó la Unidad Coordinadora de Loteos, Loteadores y Desarrolladores Inmobiliarios ``en el ámbito del entonces Ministerio de Jefatura de Gabinete de Ministros''. Once meses después ya no estaba allí.

\begin{ficha}{Resolución Nº 466 D --- B.O.\ Nº 20.858, 2 de noviembre de 2020}
Dictada el 29 de octubre de 2020 por el Secretario General de la Gobernación, con firma de Posadas.

\textbf{Artículo 1º:} encomendar a la \textbf{Secretaría de Tierras y Bienes}, dependiente de la Secretaría General de la Gobernación, las competencias asignadas a la Unidad Coordinadora de Loteos, Loteadores y Desarrolladores Inmobiliarios.

Fundamentos: la Ley 8.171 asignó a la Secretaría General de la Gobernación la competencia sobre catastro, registro de la propiedad inmobiliaria privada y fiscal y demás derechos reales (art.\ 28 inc.\ 12), y dispuso que toda norma sobre creación o funcionamiento de instituciones públicas \textbf{queda automáticamente modificada} por las competencias que ella establece (art.\ 33). Los Decretos 16/19 y 127/20 habían creado la Secretaría de Tierras y Bienes del Estado.
\end{ficha}

\begin{cautela}
Dos consecuencias, una fáctica y otra estructural.

\textbf{La fáctica identifica una foja.} La Resolución 646D/21 acredita que ``la Unidad Coordinadora de Loteos, Loteadores y Desarrolladores ha tomado la intervención que le compete'' a fojas 137, en noviembre de 2021. Un año antes de esa intervención, esas competencias habían sido encomendadas a la Secretaría de Tierras y Bienes del Estado. \textbf{La foja 137 del expediente de El Durazno es, entonces, un dictamen de la Secretaría de Tierras y Bienes del Estado}, y a esa oficina hay que pedírsela.

Eso une dos hilos que este libro venía llevando por separado. La oficina que controla los loteos privados es la misma que administra la tierra fiscal y que dictó la Resolución 82/16 sobre asignación de lotes de interés social, cuyo anexo sólo se publicó como imagen escaneada. El capítulo \ref{cap:tierrafiscal} y este capítulo tienen el mismo destinatario.

\textbf{La estructural es sobre la jerarquía del instrumento.} Un órgano de control creado por decreto, dentro del Ministerio de Jefatura de Gabinete, fue trasladado a otra dependencia por una \textbf{resolución} de la Secretaría General, sin necesidad de un nuevo decreto. El artículo 33 de la Ley 8.171 ---derogada en noviembre de 2025 por la Ley 8.511--- lo habilitaba expresamente al declarar automáticamente modificada toda norma de creación de instituciones públicas.

Este libro no objeta la legalidad de esa técnica, que es la que la propia ley previó. Registra su efecto: el único órgano específicamente creado para coordinar el control de loteos en la provincia puede cambiar de dependencia por un acto de rango inferior al que lo creó. Su ubicación institucional es más móvil que su competencia.
\end{cautela}

\subsection*{El régimen anterior, que resultó ser lo contrario de lo que este libro supuso}

Se obtuvo el texto del Decreto Nº 1.410, publicado en el Boletín Oficial Nº 9.371 del 24 de octubre de 1973. Corrige una hipótesis que este capítulo sostuvo antes de leerlo.

\begin{ficha}{Decreto 1.410/73 --- Reglamentación para la aprobación de planos de loteos urbanos y suburbanos}
\begin{itemize}[leftmargin=*,nosep]
\item \textbf{Art.\ 1.} Todo fraccionamiento con trazado de calles requiere anteproyecto \textbf{visado por la Municipalidad}, presentado en cinco copias ante la Dirección General de Inmuebles.
\item \textbf{Art.\ 3.} El plano debe consignar, entre otras cosas, \textbf{curvas de nivel con equidistancia vertical de un metro}, \textbf{zona de reserva para uso público} y \textbf{anteproyecto de obras y servicios}.
\item \textbf{Art.\ 4.} La memoria debe indicar los servicios existentes, los \textbf{servicios y mejoras que el propietario se compromete a incorporar}, y la justificación de la reserva para uso público. La firma del propietario se autentica ante escribano.
\item \textbf{Arts.\ 5 a 8.} Intervienen y \textbf{resuelven} cuatro reparticiones: Arquitectura (educación, seguridad, sanidad, abastecimiento diario, culto), Viviendas (justificación del destino y vinculación con planes oficiales), \textbf{Administración General de Aguas de Salta} (agua corriente, desagües, energía eléctrica, servicios sanitarios) y Vialidad (rutas afectadas, zona de camino, empalmes, trazado de calles).
\item \textbf{Art.\ 11.} Verificado el proyecto, se eleva al Poder Ejecutivo para su \textbf{aprobación por decreto}.
\item \textbf{Art.\ 12.} El decreto debe dejar expresa constancia de que, \textbf{previo a la venta de los lotes, el propietario debe cumplir las obligaciones contraídas}, indicando cuáles son.
\item \textbf{Art.\ 13.} Sólo dos excepciones a esa regla: loteos de planes oficiales de vivienda, o propietarios que presenten \textbf{fianza a satisfacción del Poder Ejecutivo} garantizando la ejecución de las obras en un plazo determinado.
\item \textbf{Art.\ 14.} \textbf{El cumplimiento lo verifica la Municipalidad correspondiente.}
\item \textbf{Art.\ 15.} Recién efectuada esa verificación, la Dirección General de Inmuebles \textbf{autoriza la venta por resolución expresa}, y en esa oportunidad \textbf{asigna nomenclatura catastral y número de matrícula}.
\end{itemize}
Firmado por el gobernador Ragone y su gabinete. Deroga el Decreto 8.369/70.
\end{ficha}

\begin{cautela}
El régimen de 1973 era \textbf{exigente y de secuencia estricta}. Pedía curvas de nivel metro a metro, que es información topográfica suficiente para ver una llanura de inundación. Pedía reserva para uso público y su justificación. Pedía que la Administración General de Aguas resolviera sobre agua corriente y desagües antes de aprobar. Y ataba las tres cosas que importan en un orden que no admite inversión: primero las obras, después la venta, después la matrícula.

Los considerandos lo enuncian sin ambigüedad: el fraccionamiento ``no debe consistir lisa y llanamente en la subdivisión sino que además debe contar con las obras y servicios que hagan posible el desenvolvimiento normal de la vida humana'', y por lo tanto ``el propietario que resuelva fraccionar su propiedad en lotes debe prever la ejecución de tales obras y la provisión de los servicios''.
\end{cautela}

\begin{cautela}
De donde se sigue una reformulación del argumento del libro. El problema del corredor norte no fue la ausencia de una norma exigente. \textbf{La norma exigente existió durante cuarenta y seis años.} Lo que ocurrió en El Durazno --- lotes vendidos y habitados años antes de que hubiera plano aprobado y matrícula asignada --- es precisamente la secuencia que los artículos 12, 14 y 15 del Decreto 1.410/73 prohibían. No hubo un vacío que se llenó tarde: hubo una regla que no se aplicó.

Y eso reubica al Municipio. El artículo 14 le asigna a la Municipalidad la verificación del cumplimiento de las obligaciones del loteador, y el artículo 1 le asigna el visado del anteproyecto. En el régimen anterior a 2019 el municipio no era un espectador del trámite provincial: era el verificador. El capítulo \ref{cap:hacienda} discute la capacidad municipal; este artículo indica qué tarea concreta esa capacidad debía cubrir. En el régimen de 2019 el municipio conserva el visado final, pero su silencio de veinte días se suple con una declaración jurada del desarrollador.

Corresponde una prevención que este libro no puede saltear. No está probado que el Decreto 1.410/73 estuviera vigente y fuera aplicable a cada uno de los fraccionamientos del corredor, ni cuántos de ellos tramitaron bajo él, ni si alguno usó la fianza del artículo 13. Establecer eso requiere los expedientes de aprobación de cada loteo, que este relevamiento no consultó. Lo que queda establecido es cuál era la regla escrita, y que la regla escrita era severa.
\end{cautela}

Queda además una comparación abierta que puede ser el hallazgo más productivo del capítulo. El artículo 15 del decreto de 1973 condiciona la asignación de matrícula al cumplimiento previo de las obras. El artículo 2 de la Resolución 646D/21 --- analizada más arriba en este capítulo --- condiciona la asignación de matrículas de El Durazno a dos cosas distintas: la aceptación de la donación del espacio de uso institucional y los certificados de libre deuda del período fiscal 2021.

Obras, en 1973. Donación y libre deuda, en 2021. El anexo del Decreto 1682/19, leído después, resuelve la mitad de la pregunta. No dejó de exigir obras ni certificados: los exige antes de aprobar el plano. Pero cambió la secuencia, porque admite la preventa con caución y la matrícula con obras parciales, y reemplaza la verificación municipal por un visado con plazo que una declaración jurada puede suplir. Las dos condiciones de la Resolución 646D/21 son las últimas de una lista; si las anteriores se cumplieron es lo que dirán las fojas del expediente.

\section{Dónde convergen los dos circuitos}

El artículo 11 de la reglamentación crea la \textbf{Unidad Coordinadora de Loteos, Loteadores y Desarrolladores Inmobiliarios}, que es el organismo llamado a llevar el registro y el contralor de los proyectos.

Como ya se vio, la \textbf{Resolución 466 D/20} le encomendó esas competencias a la Secretaría de Tierras y Bienes del Estado.

Éste es el hallazgo que cierra la primera tesis del libro.

\textbf{El mismo organismo que administra y adjudica la tierra fiscal a las familias de menores ingresos --- la Secretaría de Tierras y Bienes del Estado --- es también el que ejerce el registro y el contralor de los loteos privados de la provincia.}

Los dos circuitos que este libro describió como paralelos, con normas propias, ritmos propios y poblaciones destinatarias distintas, no sólo coincidieron en el organigrama del entonces Ministerio de Infraestructura, como mostrará el capítulo \ref{cap:tierrafiscal} a propósito de la Dirección General de Inmuebles. Convergieron, durante un tramo, en una misma repartición: la Dirección General de Inmuebles fue una dependencia de la Secretaría de Tierras y Bienes del Estado desde diciembre de 2019 (Decreto 127/20), el control de los loteos se sumó en octubre de 2020, y la reunión duró hasta septiembre de 2021, cuando el Decreto 844/21 pasó Inmuebles a depender directamente del ministro. Desde entonces el registro y el contralor de los loteos siguen en la Secretaría, e Inmuebles a su lado, bajo el mismo ministro y ---desde diciembre de 2025--- bajo el jefe de Gabinete (capítulo \ref{cap:tierrafiscal}).

Y la Secretaría, desde diciembre de 2023, la conduce la ex senadora provincial por el departamento La Caldera.

Dos normas complementarias completan el cuadro y conviene registrarlas.

\textbf{Decreto Nº 308/22:} delega en el Ministro de Infraestructura la facultad de aceptar las donaciones de inmuebles para uso institucional provincial que se realicen en el marco de la Ley 2.308. Es la norma que habilitó la Resolución Delegada 369 D/2022, por la cual se aceptó la donación de 8.126,79 m² sobre la matrícula 6.558 de El Durazno.

\textbf{Resolución Conjunta Nº 001/24 del ENRESP, Nº 173/24 de la Secretaría de Recursos Hídricos y Nº 683/24 de la Secretaría de Ambiente}, del 24 de septiembre de 2024: habilita un período extraordinario de regularización del \textbf{servicio de tratamiento de efluentes cloacales operado de hecho por particulares en área no servida}. Es el reconocimiento oficial de que existen privados operando cloacas sin habilitación en zonas donde la prestadora no llega, que es exactamente la situación que los vecinos plantearon en la audiencia pública de Lares de la Inmaculada en 2019.
\subsection*{Los Maitines: el caso comparado}

En San Lorenzo, municipio vecino del mismo corredor, un amparo ambiental impulsado con intervención del Ing. Pablo Campos derivó en un peritaje ordenado en 2023 por la jueza María Guadalupe Villagrán, del Juzgado Civil y Comercial de Décima Nominación, sobre el barrio privado \textit{Los Maitines S.A.}, con impugnación por la categorización del área como Categoría II bajo la Ley 26.331 y alambrado de sectores del río San Lorenzo. El desarrollador identificado en prensa es Alejandro Cornejo.

Es la única verificación parcial de H.22 que este relevamiento produjo: un apellido documentado en el cuerpo principal del libro anterior apareciendo como desarrollador inmobiliario sobre yungas. Corresponde subrayar dos límites. Es otro municipio, de modo que no acredita nada sobre La Caldera ni Vaqueros. Y proviene de prensa en cuanto a la identificación del desarrollador, aunque el peritaje y la carátula sean judiciales.
